\section{Artifact Map and Generation Contract}
\label{app:artifacts}

Every paper value is read from machine-readable artifacts by
\path{scripts/plot_paper_d.py}. The script writes white-background figures,
table fragments, the raw-frame selection receipt, and paper snapshots.
Generated paper files are self-contained, so compilation does not require the
gitignored output tree.

\begin{table}[h]
\centering\scriptsize
\caption{Protocol and artifact map. ``Official'' denotes a released host
checkpoint. The DINO-feature breadth host is trained locally and is not
official DINO-WM.}
\label{tab:artifact-map}
\begin{tabularx}{\linewidth}{p{1.1cm}p{2.2cm}Yp{2.2cm}}
\toprule
protocol & host/tasks & source & paper products\\
\midrule
Raw breadth & DINO-feature host, 10 OGBench environments &
\path{outputs/cem_raw_ogbench/report.json};
\path{outputs/cem_raw_ogbench/cells/**/result.json} &
Figs.~\ref{fig:architecture}--\ref{fig:raw-causal}\\
Official validation & released DINO-WM, Wall &
\path{outputs/cem_event_versioning_dinowm_official_v1/report.json} &
Fig.~\ref{fig:official}\\
Memory interface & released LeWM, PushT &
\path{outputs/cem_lewm_\{semantic_adapter,memory_tokens,memory_experts,adaln_memory,hybrid_interface,lora_memory\}_v1/report.json} &
Fig.~\ref{fig:lewm}\\
\bottomrule
\end{tabularx}
\end{table}

The checked-in raw snapshot is
\path{generated_results/raw_ogbench_snapshot.json}. It records the source
report hash and embeds the aggregate result used by the manuscript.
\path{generated_results/architecture_example_receipt.json} binds the selected
environment, episode, timestamps, change scores, raw cache hash, and individual
frame hashes.

\section{Raw Protocol Contract}
\label{app:contract}

\paragraph{Allowed inputs.}
The feature stage opens only `frames` and `actions` from each cache. Array
indices are timestamps. The model does not consume reward, goal, simulator
state, or any event annotation.

\paragraph{Disallowed inputs and procedures.}
The runtime source audit rejects calls named \texttt{draw\_cue},
\texttt{inject\_cue\_sequence},
\texttt{inject\_cue\_sequence\_mode}, or
\texttt{\_saliency\_map}. Cached cue-label and cue-position arrays are listed
as ignored. No fixed event interval, color rule, saturation rule, manually
selected key frame, or semantic event label is used by feature extraction,
host training, CEM training, threshold selection, or test inference.

\paragraph{Contract metadata.}
Every receipt and result stores
\texttt{cue\_window=null},
\texttt{cue\_window\_used\_by\_model=false},
\texttt{cue\_metadata\_used\_by\_model=false}, and the AST audit result. All
18 cells pass. The architecture frames are unmodified cache arrays; only
display resizing is applied by the plotting library.

\section{Data Split and Feature Host}
\label{app:raw-host}

Each environment contains 768 trajectories, 22 frames per trajectory, and 21
actions. A fixed split seed, 20,260,721, assigns 538 trajectories to train,
115 to validation, and 115 to test. Split-index SHA256 hashes are stored in
each feature receipt. The same split is used across optimization seeds.

The frozen visual encoder is DINOv2 ViT-S/14. Its local weights have SHA256
\texttt{b938bf1bc15cd2ec0feacfe3a1bb553f}\allowbreak
\texttt{e8ea9ca46a7e1d8d00217f29aef60cd9}.
The feature is a 1x1 plus 2x2 spatial pyramid over
\texttt{x\_norm\_patchtokens}. PCA is fit on train trajectories only and
retains 96 components. The finite host window contains two latent frames and
aligned actions. The residual MLP predicts the next latent and is trained only
on train trajectories. CEM training starts after the best validation host is
restored and all host parameters are frozen.

\begin{table}[h]
\centering\scriptsize
\caption{Raw host quality. MSE is held-out one-step latent error.
Host/persistence below one means the action-conditioned predictor improves on
copying the latest latent. ``Reliable'' reports passed cells.}
\label{tab:raw-host}
\begin{tabular}{lrrrr}
\toprule
environment & $n$ & host MSE & host/persistence & reliable\\
\midrule
antmaze-large & 1 & 0.2082 & 0.873 & 1/1 \\
cube-double & 1 & 0.3741 & 0.862 & 1/1 \\
cube-single & 3 & 0.3806 & 0.856 & 3/3 \\
cube-triple & 3 & 0.3519 & 0.872 & 3/3 \\
humanoidmaze-large & 1 & 0.1627 & 0.883 & 1/1 \\
pointmaze-giant & 1 & 0.2687 & 0.848 & 1/1 \\
pointmaze-large & 3 & 0.2761 & 0.866 & 3/3 \\
pointmaze-teleport & 3 & 0.3365 & 0.858 & 3/3 \\
puzzle-3x3 & 1 & 0.5694 & 0.789 & 1/1 \\
scene & 1 & 0.4544 & 0.828 & 1/1 \\
\bottomrule

\end{tabular}
\end{table}

\section{Raw CEM Training Details}
\label{app:raw-training}

\paragraph{Proposal and grouping.}
For each frame, the frozen host supplies one-step latent MSE and DINO features
supply temporal semantic change. The 0.78 train quantile sets both thresholds.
Adjacent crossings form one group. If a trajectory has no crossing, its
largest combined train-calibrated score supplies one provisional group; this
fallback is automatic and label-free.

\paragraph{Semantic keys and versions.}
A train-only MiniBatchKMeans model clusters event vectors into at most 16
semantic keys. A provisional event remains separate from the active version
with the same key. At its verification time, the new version is promoted only
if its predicted score exceeds both the validation-selected threshold and the
active score plus a hysteresis margin. A failed replacement leaves the active
version unchanged.

\paragraph{Causal-effect target.}
For candidate group $G$, the calibration target is
\[
y_G =
\frac{\mathcal L_{\mathrm{no\ memory}}-
      \mathcal L_{\mathrm{singleton}\ G}}
     {\mathcal L_{\mathrm{no\ memory}}+\epsilon}.
\]
Deleting $G$ from the singleton store returns the no-memory condition, so this
is the normalized loss increase caused by group deletion. The estimator uses
Smooth-L1 regression and within-query pairwise ranking. This target does not
use a semantic event label.

\paragraph{Validation selection.}
The validation sweep covers promotion-score quantiles, two hysteresis margins,
budgets 1/2/4, verification delays 1/3, and top-$k$ values 1/2 where legal.
The selected point minimizes future-latent MSE with a negligible retrieval
tie-break cost. Test outcomes do not enter selection. Test-only budget curves
add budget 8, and delay curves evaluate delays 1/2/3/4/6.

\section{Full Raw Controls and Telemetry}
\label{app:raw-controls}

\begin{table}[h]
\centering\scriptsize
\caption{Cell-level mean four-step test MSE over 18 cells. Every condition uses
the same frozen host, trajectories, prediction target, and memory adapter.}
\label{tab:raw-controls}
\begin{tabular}{lr}
\toprule
condition & future-latent MSE\\
\midrule
CEM & 0.48383\\
no memory & 0.48725\\
reset memory & 0.48725\\
shuffled-episode memory & 0.48939\\
matched-norm random memory & 0.48952\\
recent-only state & 0.48320\\
\bottomrule
\end{tabular}
\end{table}

The mean relative CEM improvement over no memory is $0.719\%$ with a normal
cell-level 95\% interval $[0.621,0.818]\%$. The corresponding horizon
reductions are $0.227\%$, $0.487\%$, $0.817\%$, and $1.086\%$ for horizons
one through four. All 10 environment means are positive. CEM is better than
recent-only in 3/18 cells.

Mean proposals per query are 1.923. Delayed verification promotes 0.586 and
rejects 0.414 of eligible groups. Mean active occupancy is 0.901 and retrieval
rate is 0.803. The per-cell JSON additionally records superseded, fallback,
evicted, and stale-selection counts.

\section{Causal Deletion and Action Use}
\label{app:raw-use}

For every test query with at least two candidate groups, the deletion test
compares the predicted-high singleton group with a randomly selected candidate
from the same query. Mean loss increases are 0.004253 and 0.004147. The
high-minus-random mean is 0.000106 with interval
$[-0.000056,0.000269]$ and 12/18 cell wins. CE-score Spearman correlation is
0.037 on average, positive in 13/18 cells. These values support only a weak
ranking claim.

The bounded retrieve-then-verify audit accepts 64.6\% of routed candidates
when the measured singleton group effect is positive. Its mean post-hoc
verified-memory MSE is 0.48137. This audit is computed only after observing the
future target and is not used by the primary prediction.

Action-sequence identification uses 256 test queries per cell where available.
The observed four-action sequence is placed among seven sequences sampled
from other episodes. Memory retrieval is selected with candidate actions
masked, then each action sequence is rolled out. The score is latent error to
the observed future. Mean top-1 accuracy is 0.3066 with memory and 0.3030
without memory, with 8/18 cell wins. This is an offline identification test,
not model-predictive control.

\section{Official DINO-WM Validation}
\label{app:official}

The separate official protocol uses checkpoint
\path{outputs/dinowm_wall_audit_v1/checkpoint/model_latest.pth}, epoch 65,
SHA256
\texttt{8441971becdae934fe08de5b163398390}\allowbreak
\texttt{a32f6fe1fb0a8df1}\allowbreak
\texttt{13290e2468e142b}.
The source revision is
\texttt{0a9492fa12044b852ae9e001c}\allowbreak
\texttt{c74604b79c8bb0c}. The checkpoint, action/proprioception encoders, and
DINO features remain fixed.

\begin{table}[h]
\centering\scriptsize
\caption{Released DINO-WM Wall factorial, means over seeds 0--2.
Full/control is balanced accuracy; loss is memory/reset.}
\label{tab:official-full}
\begin{tabular}{lrrrr}
\toprule
condition & overwrite & full/control & loss & high/random deletion\\
\midrule
Immediate replacement & .219 & .400/.243 & 1.311/1.397 & .085/.000\\
Hysteresis only & .751 & .774/.244 & 1.160/1.397 & .237/.000\\
Versions, no verification & .551 & .635/.236 & 1.223/1.397 & .174/.000\\
CEM & .751 & .758/.254 & 1.165/1.397 & .296/.000\\
\bottomrule
\end{tabular}
\end{table}

This protocol relocates cached genuine encoded events in latent time rather
than rendering a raw environment stream. It validates the released host and
the event-versioning mechanism, but it is not OGBench breadth, native planning,
or evidence about unmodified event discovery.

\section{LeWM Memory Interface Details}
\label{app:lewm}

All LeWM studies assert the same encoder/predictor digest before and after
training:
\texttt{5589632959b98370ad96001523025bc2}\allowbreak
\texttt{65686e82b87376d327da18cbd555f879}.
Trainable adapters are separate from the fixed host.

\begin{table}[h]
\centering\scriptsize
\caption{LeWM six-way PushT interface results, three-seed means. Intermediate
columns report post-hoc task information. Strict ratio is future-latent loss
relative to the unmodified host when objectives are comparable.}
\label{tab:lewm-full}
\begin{tabular}{lrrrrr}
\toprule
interface & memory & intermediate & host & strict ratio & result\\
\midrule
Dense residual & .919 & .169 & .156 & not comparable & interface loss\\
Memory experts & .921 & .698 & .370 & 36.79 & prediction loss\\
Semantic bottleneck & .838 & .760 & .788 & not comparable & diagnostic only\\
AdaLN & .961 & .854 & .813 & 31.08 & prediction loss\\
Memory tokens & .984 & .983 & .819 & 13.59 & prediction loss\\
Strict hybrid & .967 & .840 & .349 & 1.034 & task information lost\\
LoRA & .742 & .760 & .185 & .993 & task information lost\\
\bottomrule
\end{tabular}
\end{table}

The semantic bottleneck preserves the pairwise structure of counterfactual
latent targets (distance Pearson 0.992), but its original cue-conditioned
objective is not comparable to strict next-latent loss. Memory tokens and
AdaLN make task information available but substantially worsen prediction.
The strict hybrid and LoRA preserve prediction but do not make the binding
available at the host output. No row satisfies both requirements.

\section{Decision-Log Schema}
\label{app:logs}

Each raw decision log records the environment, optimization seed, query
timestamp, surprise stream, semantic-change stream, proposal groups, semantic
key, event timestamp, predicted effect, router score, lifecycle transitions,
promotion or rejection, supersession, eviction, fallback, retrieved event
identifiers, and the selected store configuration. It also repeats the
no-manual-event contract. Feature vectors are stored in checkpoints rather
than expanded into JSON.

The architecture receipt uses a deterministic rule: among completed test logs,
select the retrieved event with largest proposal score, then display the
preceding frame, event frame, following frame, and query frame. For the current
paper this selects HumanoidMaze-large, seed 0, episode 384, event timestamp 12,
and query timestamp 16. The four raw frame hashes and full cache hash are in
\path{generated_results/raw_architecture_receipt.json}.

\section{Reproduction}
\label{app:reproduce}

From the repository root:
\begin{verbatim}
.venv/bin/python -m pytest -q \
  scripts/test_run_cem_raw_ogbench.py

.venv/bin/python scripts/launch_cem_raw_ogbench.py \
  --smoke --gpus 1 2

.venv/bin/python scripts/launch_cem_raw_ogbench.py \
  --campaign --gpus 0 1 2

.venv/bin/python scripts/plot_cem_raw_ogbench.py
.venv/bin/python scripts/plot_paper_d.py

cd paper_d
/home/chrislin/.TinyTeX/bin/x86_64-linux/pdflatex \
  -interaction=nonstopmode -halt-on-error main.tex
/home/chrislin/.TinyTeX/bin/x86_64-linux/pdflatex \
  -interaction=nonstopmode -halt-on-error main.tex
\end{verbatim}

The launcher accepts only physical GPUs 0, 1, and 2 and rejects GPU3. The
completed launch receipt records no failed or running jobs. Generated figures,
tables, and snapshots are intended to be included with the paper so a later
PDF build does not require private output caches.

\section{Scope and Omitted Claims}
\label{app:scope}

We do not claim real-world camera generalization, robot deployment, native
DINO-WM planning, executed raw-OGBench control, unbounded recall, cross-task
transfer, or a well-calibrated causal-effect estimator. The breadth images are
unmodified simulated environment renderings. Four environments have three
optimization seeds and six have one. The recent-only comparison and weak
deletion ordering are retained because they directly limit the main result.
